\begin{table}[ht]
\centering
\begin{tabular}{lcccccc}
\hline
\textbf{Estatística} & \textbf{OBS} & \textbf{WAVERYS} & \textbf{$H_{00}$} & \textbf{$H_{01}$} & \textbf{$H_{10}$} & \textbf{$H_{11}$} \\
\hline
Média (m) & 1,97 & 2,03 & 2,22 & 2,26 & 1,96 & 1,99 \\
Máximo (m) & 8,96 & 5,99 & 6,53 & 7,18 & 6,17 & 6,93 \\
Bias (m) & 0 & 0,07 & 0,26 & 0,29 & 0 & 0,02 \\
MAE (m) & 0 & 0,23 & 0,45 & 0,48 & 0,54 & 0,56 \\
$\sqrt{MSE_{diss}}$ (m) & 0 & 0,08 & 0,26 & 0,30 & 0,05 & 0,02 \\
$\sqrt{MSE_{disp}}$ (m) & 0 & 0,33 & 0,53 & 0,54 & 0,74 & 0,77 \\
IC (adim.) & 1,00 & 0,94 & 0,82 & 0,81 & 0,69 & 0,69 \\
$D_{Pielke}$ (adim.) & 0 & 0,97 & 1,59 & 1,60 & 2,05 & 2,08 \\
Pctl 90 (m) & 2,93 & 2,94 & 3,11 & 3,17 & 2,86 & 2,94 \\
Pctl 95 (m) & 3,43 & 3,39 & 3,51 & 3,62 & 3,24 & 3,36 \\
Pctl 99 (m) & 4,43 & 4,35 & 4,34 & 4,54 & 4,20 & 4,55 \\
\hline
\end{tabular}
\caption{Estatísticas das condições de altura significativa de ondas para a boia de Itajaí comparando os dados observados, as simulações $H_{00}$, $H_{01}$, $H_{10}$, $H_{11}$ e a reanálise WAVERYS.}
\label{estats:itajai}
\end{table}